% sections/02-related-work.tex --- Аналіз досліджень.
% Citation key map (draft.md numeric -> bib key):
%   [1]=schick2023toolformer  [2]=yao2023react  [3]=qin2024toolllm
%   [4]=patil2023gorilla      [5]=ma2024sciagent [6]=mialon2023augmented
%   [7]=wang2024codeact       [8]=kim2024llmcompiler
%   [9]=nichols2024pareval    [10]=nichols2024hpccoderv2
%   [11]=xu2025llmhpc         [12]=he2025legion
%   [13]=anthropic2024mcp     [14]=hou2025mcp  [15]=luo2025mcpsmelly
%   [16]=anisimov2017parcs    [17]=bohusevych2026parcs  [18]=bohusevych2024tsp

\citsection{Analysis of related research}

\textit{Tool-augmented language-model agents.}\quad
Two main directions can be distinguished in the literature on augmenting language models with external tools. The first, proposed in Toolformer~\cite{schick2023toolformer}, showed that a model is capable of learning, without supervised examples, when and how to call external APIs such as calculators and search engines. The second, ReAct~\cite{yao2023react}, proposed a tighter integration of reasoning and action: the model produces a chain-of-thought trace, takes an action, receives an observation, and repeats. Both works established the agent-tool loop as a general paradigm. Later, the scope of tool use expanded considerably. ToolLLM by Qin et al.~\cite{qin2024toolllm} covered 16\,000 REST endpoints; Gorilla~\cite{patil2023gorilla} addressed hallucinated API signatures in ML framework calls; SciAgent~\cite{ma2024sciagent} brought the same approach to scientific domains. Mialon et al.~\cite{mialon2023augmented} survey this family as augmented language models. In all of the aforementioned systems, a tool call is stateless and local: it takes inputs, returns a result, and exits. The agent's host process remains the only site of computation.

CodeAct~\cite{wang2024codeact} generalises the action space to arbitrary Python code executed in a persistent interpreter sandbox, enabling stateful execution across turns and access to numerical libraries and file I/O. The substrate, however, remains a single interpreter, and wall-clock cost equals the script's serial execution time. LLMCompiler by Kim et al.~\cite{kim2024llmcompiler} introduces a form of parallelism by analysing a query into a dependency graph and dispatching independent calls to existing tools concurrently, reporting latency reductions of up to a factor of 3.7 on multi-hop question-answering tasks. The said parallelism, however, applies to independent calls to existing tools, not to distributed execution of user-supplied code; the system neither provisions compute resources nor compiles new code.

\textit{Language models and parallel code generation.}\quad
A separate line of research examines whether LLMs can produce correct parallel code. ParEval by Nichols et al.~\cite{nichols2024pareval} benchmarks code generation in MPI, OpenMP, CUDA, and Kokkos across 420 problems and reports a marked correctness drop relative to serial equivalents, with the gap widening as problem complexity grows. HPC-Coder-V2~\cite{nichols2024hpccoderv2} narrows this gap through fine-tuning on a synthetic parallel-code corpus, while a recent evaluation of frontier models on classical HPC kernels~\cite{xu2025llmhpc} confirms that non-trivial memory-access patterns remain a persistent source of errors. He et al.~\cite{he2025legion} use language models to optimise mappers in the Legion distributed runtime, an effective approach, but one that assumes a parallel program already exists rather than generating one.

One can argue that requiring a language model to express parallelism explicitly degrades its output quality. A model that produces a clean serial loop frequently fails to render the same loop correctly when synchronisation primitives such as \texttt{MPI\_Bcast}, \texttt{\#pragma omp parallel for}, or a CUDA kernel launch are involved.

\textit{The Model Context Protocol.}\quad
The Model Context Protocol, introduced by Anthropic in 2024~\cite{anthropic2024mcp}, specifies a standard JSON-RPC interface for agent--tool integration. The survey of Hou et al.~\cite{hou2025mcp} catalogues the early MCP ecosystem and identifies security concerns including prompt injection and tool poisoning, but reports no published work that uses MCP to expose distributed computing infrastructure. A separate analysis of public MCP tool descriptions~\cite{luo2025mcpsmelly} found that more than half of the surveyed tools provide inadequate documentation of purpose or parameters; the said finding motivated the verbosity of the tool-description design adopted in the following section.

\textit{The PARCS framework.}\quad
PARCS (Parallel Asynchronous Recursive Control Space) is a distributed computing platform whose computational model originates in early-1980s research on non-von~Neumann architectures~\cite{anisimov2017parcs}. Its core abstraction is the \textit{control space}: a finite, dynamically extensible graph of \textit{points} --- addressable locations where executable algorithmic modules operate --- connected by bidirectional \textit{channels}. Topologies such as trees, rings, and hypercubes are formed at runtime; modules communicate asynchronously and may spawn new points during execution. The platform has been implemented in numerous programming languages, including PASCAL, C, Fortran, Java, CUDA, Python, and C\#, with earlier cloud deployments targeting Amazon~EC2 and Microsoft~Azure. The present work uses its current incarnation --- PARCS-C\#, a .NET\,10 Kubernetes-native implementation that packages daemon workers as containerised pods orchestrated by KEDA and exposes job management through a REST Host~API~\cite{bohusevych2026parcs}. A previous publication~\cite{bohusevych2024tsp} extended this deployment with KEDA-driven autoscaling and evaluated it on a distributed island-model genetic algorithm for the Travelling Salesman Problem, showing 20--48\,\% shorter routes than sequential execution within the same time budget with 13--25\,\% communication overhead. That result establishes the viability of PARCS as a parallel runtime; it does not address the agent-facing interface that is the subject of this paper.

\textit{Identified gap.}\quad
Three observations follow from the literature reviewed above. First, every published agent framework executes agent-generated code within a single-threaded substrate, which limits the effective compute envelope available to the agent. Second, attempts to obtain parallelism by requesting the agent to generate parallel code conflict with the documented deficiencies of LLMs in this task. Third, no prior work employs MCP as the interface between an agent and a distributed compute runtime, despite MCP being well-suited to the role. The aim of the present work is therefore to design and describe an architectural pattern that addresses all three gaps simultaneously.
